%%%%%%%%%%%%%%%%%%%%%%%%%%%%%%%%%%%%%%%%%%%%%%%%%%%%%%%%%%%%%%%%%%%%%%
% Lớp 11 | Giới hạn. Hàm số liên tục | Giới hạn của dãy số
% Dạng 1D3?1-1: Câu hỏi lý thuyết
% Số câu: 40  (sinh tự động bởi tach_cau.py; từ nay thầy sửa trực tiếp file này)
%%%%%%%%%%%%%%%%%%%%%%%%%%%%%%%%%%%%%%%%%%%%%%%%%%%%%%%%%%%%%%%%%%%%%%

%%%%%%%%%%%%%%%% TRẮC NGHIỆM 4 PHƯƠNG ÁN %%%%%%%%%%%%%%%%

% nguồn: Ok_Toan11_TracNghiem_Lan2
\begin{ex}%[1D3N1-1]%[ma:61d43498]
	Cho dãy số $\left(u_n\right)$ thỏa mãn  $\lim u_n=\dfrac{3}{2}$. Giá trị của $\lim\left(u_n+4\right)$ bằng
	 \choice
		 {\True $\dfrac{11}{2}$}
		{$\dfrac{11}{4}$} 
		{$\dfrac{13}{2}$} 
		{$\dfrac{13}{4}$} 
	\loigiai{ 
		Ta có $\lim\left(u_n+4\right)=\dfrac{3}{2}+4=\dfrac{11}{2}$. 
	} 
\end{ex}

% nguồn: Ok_Toan11_TracNghiem_Lan2
\begin{ex}%[1D3N1-1]%[ma:524db969]
	 Cho $\lim u_n=-3$. $\lim v_n=2$. Khi đó $\lim\left(u_n-v_n\right)$ bằng
	 \choice
		 {\True $-5$}
		 {$-1$}
		 {$5$}
		 {$1$}
	\loigiai{Ta có  $\lim\left(u_n-v_n\right)=\lim u_n-\lim v_n=-3-2=-5$.
	} 
\end{ex}

% nguồn: Ok_Toan11_TracNghiem_Lan2
\begin{ex}%[1D3N1-1]%[ma:60b2a08c]
	Cho dãy số $\left(u_n\right)$ thỏa  mãn $\lim u_n=-5$. Giá trị của $\lim\left(u_n-2\right)$ bằng
	 \choice
		{$3$} 
		{\True  $-7$}
		{$10$} 
		 {$-10$}
	\loigiai{ 
		Ta có $\lim\left(u_n-2\right)=(-5)+(-2)=-7$. 
	 }
\end{ex}

% nguồn: Ok_Toan11_TracNghiem_Lan2
\begin{ex}%[1D3N1-1]%[ma:ea4c22e7]
	Cho  dãy số $\left(u_n\right)$, $\left(v_n\right)$ thỏa mãn $\lim u_n=11$, $\lim v_n=4$. Giá trị của $\lim\left(u_n+v_n\right)$ bằng
	 \choice
		 {$4$}
		{$7$} 
		{$11$} 
		{\True $15$} 
	\loigiai{ 
		Ta  có $\lim\left(u_n+v_n\right)=\lim u_n+\lim v_n=11+4=15$.
	} 
\end{ex}

% nguồn: Ok_Toan11_TracNghiem_Lan2
\begin{ex}%[1D3N1-1]%[ma:5aff3021]
	Phát biểu nào  sau đây là {\bf sai}?
	 \choice
		{\True $\lim  q^n=0$, ($|q|>1$)}
		{$\lim C=C$, ($C$ là hằng  số)}
		{$\lim\dfrac{1}{n^k}=0$,  ($k\in\mathbb{N}^*)$}
		 {$\lim\dfrac{1}{n}=0$}
	 \loigiai{
	Vì $\lim q^n=0$ nếu $|q|<1$.} 
\end{ex}

% nguồn: Ok_Toan11_TracNghiem_Lan2
\begin{ex}%[1D3N1-1]%[ma:04321f11]
	Cho dãy $\left(u_n\right)$ có $\lim u_n=3$, dãy $\left(v_n\right)$ có $\lim v_n=5$. Khi đó $\lim\left(u_nv_n\right)$  là
	\choice 
		{\True  $15$}
		 {$8$}
		 {$5$}
		 {$3$}
	 \loigiai{
		Ta có $\lim\left(u_n \cdot v_n\right)=\lim u_n\cdot \lim  v_n=3\cdot 5=15$.
	} 
\end{ex}

% nguồn: Ok_Toan11_TracNghiem_Lan2
\begin{ex}%[1D3N1-1]%[ma:8112a894]
	Cho các  dãy số $\left(u_n\right)$, $\left(v_n\right)$ và $\lim{u_n}=a$, $\lim{v_n}=+\infty$ trong đó $a$ hữu hạn thì $\lim\dfrac{u_n}{v_n}$ bằng
	\choice 
		 {$1$}
		{\True  $0$}
		 {$-\infty $}
		{$+\infty $} 
	\loigiai{ 
		Theo tính chất giới hạn ta có \lq\lq cho dãy số $\left(u_n\right)$, $\left(v_n\right)$ và $\lim{u_n}=a$, $\lim{v_n}=+\infty $ trong  đó $a$ hữu hạn thì $\lim\dfrac{u_n}{v_n}=0$\rq\rq.
	} 
\end{ex}

% nguồn: Ok_Toan11_TracNghiem_Lan2
\begin{ex}%[1D3N1-1]%[ma:51ad4b01]
	Cho hai dãy số $\left(u_n\right),$ $\left(v_n\right)$ thỏa mãn $\lim u_n=-4$ và $\lim v_n=3$. Giá trị của $\lim\left(u_n  \cdot v_n\right)$ bằng
	 \choice
		 {$12$}
		 {\True $-12$}
		{$1$} 
		{$7$} 
	 \loigiai{
		Ta có $\lim\left(u_n \cdot  v_n\right)=\lim u_n \cdot \lim v_n=(-4) \cdot 3=-12$.
	} 
\end{ex}

% nguồn: Ok_Toan11_TracNghiem_Lan2
\begin{ex}%[1D3N1-1]%[ma:fdcdceb3]
	Cho hai dãy số $\left(u_n\right)$ và  $\left(v_n\right)$ có $u_n=\dfrac{1}{n+1}$; $v_n=\dfrac{3}{n+3}$. Tính $\lim\dfrac{u_n}{v_n}$.
	 \choice
		 {$0$}
		 {$ 3$}
		{\True  $\dfrac{1}{3}$}
		 {$+\infty $}
	\loigiai{ 
	Ta  có $I=\lim\dfrac{u_n}{v_n}=\lim\dfrac{\dfrac{1}{n+1}}{\dfrac{3}{n+3}}=\lim\dfrac{n+3}{3\left(n+1\right)}=\lim\dfrac{1+\dfrac{3}{n}}{3\left(1+\dfrac{1}{n}\right)}=\dfrac{1+0}{3\left(1+0\right)}=\dfrac{1}{3}$.}
\end{ex}

% nguồn: Ok_Toan11_TracNghiem_Lan2
\begin{ex}%[1D3N1-1]%[ma:737358f5]
	Trong các khẳng định dưới đây có bao nhiêu  khẳng định đúng?
	\begin{itemize}
		\item $\lim{n^k}=+\infty $ với $k$ nguyên dương.  $\qquad (I)$
		\item $\lim{q^n}=+\infty $ nếu $\left| q\right|<1$. $\qquad (II)$ 
		\item $\lim{q^n}=+\infty  $ nếu $q>1$. $\qquad (III)$
	\end{itemize}
	\choice 
		 {$0$}
		{$1$} 
		 {$3$}
		 {\True $2$}
	\loigiai{ 
		\begin{itemize}
			\item $\lim{n^k}=+\infty $ với $k$  nguyên dương $\Rightarrow(I)$ là khẳng định đúng.
			\item $\lim{q^n}=+\infty $ nếu $\left| q\right|<1$ $\Rightarrow\left(II\right)$ là khẳng định sai vì $\lim{q^n}=0$ nếu  $\left| q\right|<1$.
			\item $\lim{q^n}=+\infty $ nếu $q>1$  $\Rightarrow\left(III\right)$ là khẳng định đúng.
		\end{itemize}
	Vậy số khẳng định đúng  là $2$.}
\end{ex}

% nguồn: Ok_Toan11_TracNghiem_Lan2
\begin{ex}%[1D3N1-1]%[ma:0f980d67]
	Nếu $\lim  u_n=-3 $; $\lim v_n=1$ thì $\lim\left(u_n+v_n\right)$ bằng
	\choice 
		{$-1$} 
		{$1$} 
		 {\True $-2$}
		 {$-4$}
	\loigiai{ 
		 Ta có $\lim\left(u_n+v_n\right)=\lim u_n+\lim v_n=-3+1=-2$.
	} 
\end{ex}

% nguồn: Ok_Toan11_TracNghiem_Lan2
\begin{ex}%[1D3N1-1]%[ma:390f8598]
	Cho dãy số $\left(u_n\right)$ thỏa mãn $\lim\left(u_n-2\right)=0$. Giá trị của $\lim u_n$ bằng 
	 \choice
		{\True $2$} 
		 {$-2$}
		 {$1$}
		{$0$} 
	\loigiai{ 
		Ta có $\lim\left(u_n-2\right)=0 \Leftrightarrow \lim u_n-\lim 2=0\Leftrightarrow  \lim u_n=2$.
	 }
\end{ex}

% nguồn: Ok_Toan11_TracNghiem_Lan2
\begin{ex}%[1D3N1-1]%[ma:88f7e635]
	 Cho hai dãy số $\left(u_n\right)$ và $\left(v_n\right)$ thoả mãn $\lim u_n=6$ và $\lim v_n=2$. Giá trị của $\lim\left(u_n-v_n\right)$ bằng
	\choice 
		 {$12$}
		{$8$} 
		 {$-4$}
		{\True  $4$}
	 \loigiai{
		 Ta có $\lim\left(u_n-v_n\right)=\lim u_n-\lim v_n=6-2=4$.
	} 
\end{ex}

% nguồn: Ok_Toan11_TracNghiem_Lan2
\begin{ex}%[1D3N1-1]%[ma:cb8d1b13]
	Phát biểu  nào sau đây là \textbf{sai}?
	\choice 
		{$\lim{u_n}=c$ ($u_n=c$ là hằng  số)}
		{\True $\lim{q^n}=0$ $\left(\left| q\right|>1\right)$} 
		 {$\lim\dfrac{1}{n}=0$}
		{$\lim\dfrac{1}{n^k}=0$ $\left(k>1\right)$} 
	\loigiai{ 
	Theo định nghĩa giới hạn hữu hạn của  dãy số thì $\lim{q^n}=0$ $\left(\left| q\right|<1\right)$.}
\end{ex}

% nguồn: Ok_Toan11_TracNghiem_Lan2
\begin{ex}%[1D3N1-1]%[ma:67a82e8d]
	Cho hai dãy số $\left(u_n\right)$, $\left(v_n\right)$ thỏa mãn $\lim u_n=2$, $\lim v_n=-3$.  Giá trị của $\lim\left(u_n \cdot v_n\right)$ bằng
	 \choice
		 {$6$}
		 {$5$}
		{\True  $-6$}
		 {$-1$}
	\loigiai{ 
		Ta có $\lim\left(u_n \cdot  v_n\right)=2 \cdot(-3)=-6$.
	 }
\end{ex}

% nguồn: Ok_Toan11_TracNghiem_Lan2
\begin{ex}%[1D3N1-1]%[ma:9034d7fc]
	Cho dãy số $\left(u_n\right)$ thỏa $\left|u_n-2\right|<\dfrac{1}{n^3}$ với mọi  $n\in\mathbb{N}^*$ . Khi đó
	 \choice
		{$\lim{u_n}$ không tồn tại} 
		{$\lim{u_n}=1$} 
		 {$\lim{u_n}=0$}
		 {\True $\lim{u_n}=2$}
	\loigiai{Ta có $\left|u_n-2\right|<\dfrac{1}{n^3} \Rightarrow\lim\left(u_n-2\right)=\lim\dfrac{1}{n^3}=0 \Rightarrow\lim{u_n}-2=0\Rightarrow\lim{u_n}=2$.} 
\end{ex}

% nguồn: Ok_Toan11_TracNghiem_Lan2
\begin{ex}%[1D3N1-1]%[ma:9a8faaa7]
	Cho dãy $(u_n)$ có $\lim u_n=3$, dãy $(v_n)$ có $\lim v_n=5$. Khi đó  $\lim\left(u_n\cdot v_n\right)=$?
	 \choice
		 {\True $15$}
		 {$8$}
		 {$5$}
		 {$3$}
	\loigiai{Nếu $\lim u_n=a$,  $\lim v_n=b$ thì $\lim\left(u_n\cdot v_n\right)=ab$.\\
		 Do đó $\lim\left(u_n\cdot v_n\right)=3\cdot 5=15$.
	} 
\end{ex}

% nguồn: Ok_Toan11_TracNghiem_Lan2
\begin{ex}%[1D3N1-1]%[ma:6a2e47d2]
	Cho $\lim a_n=-3$, $\lim b_n=5$. Khi đó  $\lim\left(a_n-b_n\right)$ bằng
	\choice 
		{$-2$} 
		{$8$} 
		 {$2$}
		{\True $-8$} 
	 \loigiai{
		Ta có $\lim\left(a_n-b_n\right)=\lim a_n-\lim  b_n=(-3)-5=-8$.
	 }
\end{ex}

% nguồn: Ok_Toan11_TracNghiem_Lan2
\begin{ex}%[1D3N1-1]%[ma:b21ac56c]
	Cho dãy số $\left(u_n\right)$ thỏa mãn  $\lim\left(u_n+3\right)=0$. Giá trị của $\lim u_n$ bằng
	 \choice
		{$3$} 
		{\True $-3$} 
		 {$2$}
		 {$0$}
	 \loigiai{
		Ta có $\lim\left(u_n+3\right)=0 \Leftrightarrow \lim  u_n=-3$.
	 }
\end{ex}

% nguồn: Ok_Toan11_TracNghiem_Lan2
\begin{ex}%[1D3N1-1]%[ma:bf4f91b0]
	Trong các  mệnh đề dưới đây, mệnh đề nào \textbf{sai}?
	\choice 
		{Nếu $\lim u_n=+\infty$ và $\lim v_n=a>0$ thì  $\lim\left(u_n v_n\right)=+\infty$}
		{ Nếu $\lim u_n=a \neq 0$ và $\lim v_n=  \pm \infty$ thì $\lim\left(\dfrac{u_n}{v_n}\right)=0$}
		{\True Nếu $\lim u_n=a>0$ và $\lim v_n=0$ thì $\lim\left(\dfrac{u_n}{v_n}\right)=+\infty$} 
		{ Nếu $\lim u_n=a<0$ và $\lim v_n=0$  và $v_n>0$ với mọi $n$ thì $\lim\left(\dfrac{u_n}{v_n}\right)=-\infty$}
	 \loigiai{
		 Nếu $\lim u_n=a>0$ và $\lim v_n=0$ thì $\lim\left(\dfrac{u_n}{v_n}\right)=+\infty$ là mệnh đề sai vì chưa rõ dấu của $v_n$ là dương hay âm.
	} 
\end{ex}

% nguồn: Ok_Toan11_TracNghiem_Lan2
\begin{ex}%[1D3N1-1]%[ma:25ea7a60]
	Khẳng định nào sau đây  là đúng?
	\choice 
		{\True Ta nói dãy số $\left(u_n\right)$ có giới hạn là số $a$ khi $n\to+\infty $ ,  nếu $\underset{n\to+\infty}{\lim}\,\left(u_n-a\right)=0$}
		{Ta nói dãy số $\left(u_n\right)$ có giới hạn là  $0$ khi $n$ dần tới vô cực, nếu $\left|u_n\right|$ có thể lớn hơn một số dương tùy ý, kể từ một số hạng nào đó trở đi}
		{Ta nói dãy số $\left(u_n\right)$ có giới hạn $+\infty $ khi $n\to+\infty $ nếu $u_n$ có thể nhỏ hơn một số dương bất kì, kể từ một số hạng nào đó  trở đi}
		{Ta nói dãy số $\left(u_n\right)$ có giới hạn $-\infty $ khi $n\to+\infty $ nếu $u_n$ có thể lớn hơn  một số dương bất kì, kể từ một số hạng nào đó trở đi}
	\loigiai{ 
	Theo lý thyết ta có dãy số  $\left(u_n\right)$ có giới hạn là số $a$ khi $n\to+\infty $ , nếu $\underset{n\to+\infty}{\lim}\,\left(u_n-a\right)=0$. }
\end{ex}

% nguồn: Ok_TracNghiemLop11_Lan1
\begin{ex}%[1D3N1-1]%[ma:9aa15cc5]
	Cho hai dãy số $\left( u_n\right)$ và $\left( v_n\right)$  có $\lim
	u_n =6$ và  $\lim
	v_n =-10$. Tính giới hạn  $\lim
	(u_n+v_n)$. 
	\choice 
		{  $16$}
		 { $-60$}
		{$-16$} 
		{ \True $-4$} 
	\loigiai{ 
		Ta  có $\lim \left( u_n+v_n\right) =\lim u_n-\lim v_n=6-10=-4$.
	 }
\end{ex}

% nguồn: Ok_TracNghiemLop11_Lan1
\begin{ex}%[1D3N1-1]%[ma:4672c6eb]
	Cho dãy số $\left( u_n\right)$ thỏa mãn $\lim\limits_{n\to +\infty} \left(u_n+3\right)=0$. Giá trị $\lim\limits_{n\to +\infty}  u_n$ bằng
	\choice 
		 {$3$}
		 {\True $-3$}
		{$2$} 
		{$0$} 
	 \loigiai{
		Ta có $\lim\limits_{n\to +\infty} \left(u_n- 3\right)=0 \Leftrightarrow \lim\limits_{n\to +\infty} u_n=-3$. 
	 }
\end{ex}

% nguồn: Ok_TracNghiemLop11_Lan1
\begin{ex}%[1D3N1-1]%[ma:f2f5a187]
	Trong các mệnh đề sau, mệnh đề nào là mệnh đề  \textbf{đúng}?
	\choice 
		 {$\lim\limits_{x\to+\infty} \dfrac{1}{n}=1$}
		{$\lim\limits_{x\to+\infty} q^n=0,  q>1$}
		 {$\lim\limits_{x\to+\infty} 1^n=0$}
		{\True  $\lim\limits_{x\to+\infty} q^n=0,|q|<1$}
	 \loigiai{Dựa vào một số giới hạn đặc biệt:
		$\lim\limits_{x\to+\infty}  \dfrac{1}{n}=0; \lim\limits_{x\to+\infty} q^n=0 , |q|<1$.
	} 
\end{ex}

% nguồn: Ok_TracNghiemLop11_Lan1
\begin{ex}%[1D3N1-1]%[ma:bed0c85a]
	Giá trị của $\lim q^n$ với $\left|q\right|<1$  bằng
	\choice 
		 {$1$}
		{\True $0$} 
		{$+\infty$} 
		 {$-\infty$}
	\loigiai{ 
		Với  $\left|q\right|<1$ thì $\lim q^n=0$.
	 }
\end{ex}

% nguồn: Ok_TracNghiemLop11_Lan1
\begin{ex}%[1D3N1-1]%[ma:de7fb660]
	Trong các khẳng định dưới đây, khẳng định nào \textbf{sai}? 
	 \choice
		 {\True $\lim\limits_{n \to +\infty} q^n=0$}
		{$\lim\limits_{n \to +\infty} \dfrac{1}{n^k}=0$ với $k$  nguyên dương}
		 {$\lim\limits_{n \to +\infty} \dfrac{1}{n}=0$}
		{Nếu $u_n=c$ ($c$ là hằng số) thì $\lim\limits_{n \to +\infty} u_n=\lim c=c$} 
	\loigiai{ 
	Ta có $\lim\limits_{n \to +\infty} q^n=0$ khi  $|q|<1$; $\lim\limits_{n \to +\infty} q^n=+\infty$ khi $q>1$.}
\end{ex}

% nguồn: Ok_TracNghiemLop11_Lan1
\begin{ex}%[1D3N1-1]%[ma:051b4635]
	Đặt $\lim u_n  = a$, $\lim v_n = b$. Mệnh đề nào dưới đây là \textbf{sai}?
	 \choice
		{\True $\lim  ( u_n \cdot v_n)$ = $\lim u_n + \lim v_n$}
		 {$\lim (u_n + v_n)$ = $\lim u_n + \lim v_n$}
		{$\lim (u_n - v_n)$ = $\lim u_n  - \lim v_n$}
		{$\lim  (u_n \cdot v_n)$ = $\lim u_n \cdot \lim v_n$}
	\loigiai{ 
		Mệnh đề $\lim (u_n \cdot v_n)$ = $\lim u_n \cdot \lim v_n$ là mệnh đề đúng nên mệnh đề \break $\lim ( u_n \cdot v_n)$ =  $\lim u_n + \lim v_n$ là sai.
	 }
\end{ex}

% nguồn: Ok_TracNghiemLop11_Lan1
\begin{ex}%[1D3H1-1]%[ma:d72ee010]
	Phát biểu nào sau  đây \textbf{sai}?
	\choice 
		{$\lim\limits_{n  \to +\infty} \dfrac{4}{\left(\sqrt{3}\right)^n}=0$}
		{$\lim\limits_{n \to  +\infty} \left(\dfrac{2}{3}\right)^n=0$}
		{$\lim\limits_{n  \to +\infty} \left(-\dfrac{\sqrt{2}}{2}\right)^n=0$}
		{\True $\lim\limits_{n \to  +\infty} \left(\dfrac{2}{\sqrt{3}}\right)^n=0$}
	 \loigiai{
		Ta có $\dfrac{2}{\sqrt{3}}>1$ nên $\lim\limits_{n  \to +\infty} \left(\dfrac{2}{\sqrt{3}}\right)^n=+\infty$.
	} 
\end{ex}


%%%%%%%%%%%%%%%% TRẢ LỜI NGẮN %%%%%%%%%%%%%%%%

% nguồn: Ok_TraLoiNgan_Lop11_Lan1
\begin{ex}%[1D3V1-1]%[ma:96707518]
	Tìm giới hạn  sau: ${\lim \left[\dfrac{1}{1\cdot 3}+\dfrac{1}{3\cdot 5}+\dfrac{1}{5\cdot 7}+\cdots+\dfrac{1}{(2n-1)\cdot(2n+1)}\right]}$.
	\shortans{0,5} 
	 \loigiai{
		${\begin{aligned}
		& \lim \left[\dfrac{1}{1 \cdot 3}+\dfrac{1}{3 \cdot  5}+\dfrac{1}{5 \cdot 7}+\cdots+\dfrac{1}{(2 n-1)(2 n+1)}\right] \\
		 &=\lim \dfrac{1}{2}\left(1-\dfrac{1}{3}+\dfrac{1}{3}-\dfrac{1}{5}+\cdots+\dfrac{1}{2 n-1}-\dfrac{1}{2 n+1}\right)=\lim \dfrac{1}{2}\left(1-\dfrac{1}{2 n+1}\right)=\dfrac{1}{2}.
	\end{aligned}}$
} 
\end{ex}


%%%%%%%%%%%%%%%% TỰ LUẬN %%%%%%%%%%%%%%%%

% nguồn: Ok_TuanLuan_Toan11_Lan2
\begin{ex}%[1D3H1-1]%[ma:27816e72]
	Dùng định nghĩa dãy số có  giới hạn 0 tìm $\lim u_n$ với $u_n=\dfrac{{{(-1)}^n}}{3n+2}$.
	\loigiai{ 
		Với $a  > 0$ nhỏ tùy ý, ta có $\left|u_n \right|=\left|\dfrac{(-1)^n}{3n+2} \right|=\dfrac{1}{3n+2} < \dfrac{1}{3n} < a\Leftrightarrow n > \dfrac{1}{3a}$.\\
		Chọn $n_0=\left[\dfrac{1}{3a} \right]$. Do đó $\forall a > 0$, $\exists  n_0 \colon n > n_o$ ta luôn có $\left|u_n \right|< a
	\Rightarrow \lim \dfrac{(-1)^n}{3n+2}=0$. 
	} 
\end{ex}

% nguồn: Ok_TuanLuan_Toan11_Lan2
\begin{ex}%[1D3H1-1]%[ma:c3aa29fc]
	Dùng  định nghĩa dãy số có giới hạn $0$ tìm $\lim u_n$ với $u_n=\dfrac{1+\cos n^3}{2 n+3}$.
	\loigiai{ 
		Ta luôn có $\left|\cos n^3\right| \leq 1$, $\forall n \in  \mathbb{N}^*$.\\
		Với $a>0$ nhỏ tùy ý, $\left|u_n\right|=\left|\dfrac{1+\cos n^3}{2  n+3}\right| \leq\left|\dfrac{2}{2 n+3}\right| \leq\left|\dfrac{2}{2 n}\right|=\dfrac{1}{n}<a \Leftrightarrow n>\dfrac{1}{a}$.\\
		 Chọn $n_0=\left[\dfrac{1}{a}\right]$.\\
		Do đó $\forall a>0$, $\exists n_0 \colon n>n_o$ ta luôn  có $\left|u_n\right|<a \Rightarrow \lim \dfrac{1+\cos n^3}{2 n+3}=0$.
	 }
\end{ex}

% nguồn: Ok_TuanLuan_Toan11_Lan2
\begin{ex}%[1D3H1-1]%[ma:0dab41df]
	 Cho dãy số $\left(u_n \right)$ với $u_n=\dfrac{1 \cdot 3 \cdot 5 \cdot 7 \cdots (2n-1)}{2 \cdot 4 \cdot 6 \cdots 2n}$. Tính $\lim u_n$.
	 \loigiai{
		Ta có $u_n > 0$, $\forall n\in \mathbb{N}^*$, do đó $\left(u_n \right)^2 > 0$, $\forall  n\in \mathbb{N}^*$.\\
		$\Rightarrow \left(u_n \right)^2 = \dfrac{1^2 \cdot 3^2 \cdot 5^2 \cdot 7^2  \cdots (2n-1)^2}{2^2 \cdot 4^2 \cdot 6^2 \cdots (2n)^2} < \dfrac{1^2 \cdot 3^2 \cdot 5^2 \cdot 7^2 \cdots (2n-1)^2}{(2^2-1)(4^2-1)(6^2-1) \cdots ((2n)^2 - 1)}$.\\
		 Do đó
		\[ 
		\left(u_n \right)^2 < \dfrac{1^2  \cdot 3^2 \cdot 5^2 \cdots (2n-1)^2}{1 \cdot 3 \cdot 3 \cdot 5 \cdots (2n-1)(2n+1)} = \dfrac{1}{2n+1}.
		\] 
		Vậy,  với mọi $n\in \mathbb{N}^*$ thì $0 < \left(u_n \right)^2 < \dfrac{1}{2n+1}$.\\
		Mà $\lim \dfrac{1}{2n+1}  = 0$ nên $\lim \left(u_n \right)^2 = 0$.\\
		Từ đó suy ra $\lim u_n =  0$.
	} 
\end{ex}

% nguồn: Ok_TuanLuan_Toan11_Lan2
\begin{ex}%[1D3V1-1]%[ma:4217a4ab]
	Cho dãy  số $\left(u_n \right)$ với $u_n=\dfrac{\sqrt{n+1}}{n+2}$. Tính $\lim u_n$.
	\loigiai{ 
	Ta có $\left|u_n \right|=\left|\dfrac{\sqrt{n+1}}{n+2} \right|< \dfrac{\sqrt{n+1}}{n+1}=\dfrac{1}{\sqrt{n+1}} < \dfrac{1}{\sqrt{n}}$, $\forall n\in \mathbb{N}$. Vì $\lim \dfrac{1}{\sqrt{n}}=0$ nên $\lim  u_n=0$.}
\end{ex}

% nguồn: Ok_TuanLuan_Toan11_Lan2
\begin{ex}%[1D3V1-1]%[ma:93e884df]
	Chứng minh rằng  $\lim \dfrac{\sin^2n}{n+2}=0$.
	 \loigiai{
		Với $ a>0$ nhỏ tùy ý, ta có $\left| \dfrac{\sin^2n}{n+2}  \right|=\dfrac{\sin^2n}{n+2}<\dfrac{1}{n+2}<a\Leftrightarrow n>\dfrac{1}{a}-2$.\\
		 Chọn $n_0=\left[\dfrac{1}{a}-2 \right]$. Khi đó $\forall a>0$, $\exists n_0:n>n_0$ ta luôn có $\left| \dfrac{\sin^2n}{n+2} \right|<a$ $\Rightarrow \lim \dfrac{\sin^2n}{n+2}=0$.\\
	\textit{Chú ý}: Kí hiệu  $\left[a \right]$ là lấy phần nguyên của $ a$.}
\end{ex}

% nguồn: Ok_TuanLuan_Toan11_Lan2
\begin{ex}%[1D3V1-1]%[ma:fc277136]
	 Cho dãy số $\left(u_n \right)$ với $u_n=(-0{,}97)^n$. Tính $\lim u_n$.
	 \loigiai{
		Theo công thức giới hạn đặc biệt, ta có $\left|-0{,}97 \right|<  1$ nên $\lim u_n=0$.
	} 
\end{ex}

% nguồn: Ok_TuanLuan_Toan11_Lan2
\begin{ex}%[1D3V1-1]%[ma:d9546b1f]
	Cho dãy số $\left(u_n \right)$ với  $u_n=\dfrac{n+2\sin (n+1)}{n\sqrt[3]{n}+2\sqrt[3]{n}}$. Tính $\lim u_n$.
	\loigiai{ 
	Ta có $\left|u_n \right|=\left|\dfrac{n+2\sin (n+1)}{n\sqrt[3]{n}+2\sqrt[3]{n}}  \right|< \dfrac{n+2}{\sqrt[3]{n}(n+2)} < \dfrac{1}{\sqrt[3]{n}}$, $\forall n\in \mathbb{N}$. Vì $\lim \dfrac{1}{\sqrt[3]{n}}=0$ nên $\lim u_n=0$.}
\end{ex}

% nguồn: Ok_TuanLuan_Toan11_Lan2
\begin{ex}%[1D3V1-1]%[ma:e4b9908c]
	Chứng minh rằng $\lim  \dfrac{1}{n^2+1}=0$.
	 \loigiai{
		Với $ a>0$ nhỏ tùy ý,  ta có $\left|\dfrac{1}{n^2+1} \right|=\dfrac{1}{n^2+1}<a\Leftrightarrow n>\sqrt{\dfrac{1}{a}-1}$.\\
		Chọn $n_0=\left[\sqrt{\dfrac{1}{a}-1} \right]$. Khi đó $\forall a>0$, $\exists n_0:n>n_0$ ta luôn có $\left| \dfrac{1}{n^2+1}  \right|<a$ $\Rightarrow \lim \dfrac{1}{n^2+1}=0$.\\
	\textit{Chú ý}: Kí hiệu $\left[a \right]$ là lấy phần nguyên của  $ a$.}
\end{ex}

% nguồn: Ok_TuanLuan_Toan11_Lan2
\begin{ex}%[1D3V1-1]%[ma:eaffd6d3]
	Chứng minh rằng $\lim  \left( \dfrac{(-1)^n}{2^{n+1}}-\dfrac{1}{3^{n+1}} \right)=0$.
	\loigiai{ 
		Bằng phương pháp quy nạp ta chứng minh được $ n<2^n$, $\forall n\in \mathbb{N}^*\Rightarrow \dfrac{1}{2^n}<\dfrac{1}{n}$, $\forall n\in  \mathbb{N}^*$.\\
		Với  $ a>0$ nhỏ tùy ý, ta có
		$$\left|  \dfrac{(-1)^n}{2^{n+1}}-\dfrac{1}{3^{n+1}} \right|=\dfrac{1}{2^{n+1}}+\dfrac{1}{3^{n+1}}<\dfrac{1}{2^{n+1}}+\dfrac{1}{2^{n+1}}=\dfrac{1}{2^n}<\dfrac{1}{n}<a\Leftrightarrow n>\dfrac{1}{a}.$$
		Chọn  $n_0=\left[\dfrac{1}{a} \right]$. Khi đó $\forall a>0$, $\exists n_0:n>n_0$ ta luôn có
		$$\left| \dfrac{(-1)^n}{2^{n+1}}-\dfrac{1}{3^{n+1}} \right|<a\Rightarrow \lim \left( \dfrac{(-1)^n}{2^{n+1}}-\dfrac{1}{3^{n+1}}  \right)=0.$$
		\textit{Chú ý}: Kí hiệu $\left[a \right]$ là lấy phần  nguyên của $ a$.
	 }
\end{ex}

% nguồn: Ok_TuanLuan_Toan11_Lan2
\begin{ex}%[1D3V1-1]%[ma:3882b8f4]
	Cho  dãy số $\left(u_n\right)$ với $u_n=\sqrt{n^2+1}-n$. Tính $\lim u_n$.
	\loigiai{ 
		Ta có $u_n=\sqrt{n^2+1}-n=\dfrac{\left(\sqrt{n^2+1}-n \right)\left(\sqrt{n^2+1}+n \right)}{\sqrt{n^2+1}+n}=\dfrac{1}{n\left(\sqrt{1+\dfrac{1}{n^2}}+1 \right)}=\dfrac{1}{n} \cdot  \dfrac{1}{\sqrt{1+\dfrac{1}{n^2}}+1}$
	Vì $\lim \dfrac{1}{n}=0$  và $\lim \dfrac{1}{\sqrt{1+\dfrac{1}{n^2}}+1}=\dfrac{1}{2}$ nên $\lim u_n=0$. }
\end{ex}

% nguồn: Ok_TuanLuan_Toan11_Lan2
\begin{ex}%[1D3V1-1]%[ma:892fc0a1]
	Dùng định nghĩa dãy số có giới hạn 0 tìm $\lim  u_n$ với $u_n=\dfrac{\sqrt[\text{n}]{n!}}{\sqrt{n^3+2n}}$.
	\loigiai{ 
		Với $a  > 0$ nhỏ tùy ý, ta có $\left|u_n \right|=\left|\dfrac{\sqrt[\text{n}]{n!}}{\sqrt{n^3+2n}} \right|< \dfrac{\sqrt[\text{n}]n^n}{\sqrt{n^3+2n}}=\dfrac{n}{\sqrt{n^3+2n}} < \dfrac{n}{n\sqrt{n}}=\dfrac{1}{\sqrt{n}} < a\Leftrightarrow n > \dfrac{1}a^2$. Chọn $n_0=\left[\dfrac{1}{a^2} \right]$. Do đó $\forall a > 0$, $\exists n_0:n > n_o$ ta luôn có $\left|u_n \right|< a
	 \Rightarrow \lim \dfrac{\sqrt[\text{n}]{n!}}{\sqrt{n^3+2n}}=0$. }
\end{ex}
